\section{Spin and Fermions}

The coin carries spin, matter is fermionic, and the rule makes matter stable, all without adding anything to the
base. The three facts come out of the substrate's own symmetry and the measured behavior of its collision, not
from any extra ingredient grafted on for the purpose. This section reads them off in order.

Two terms recur throughout, so a plain statement of each helps. A \emph{spinor} is the mathematical object that
describes a particle with half-integer spin. Its signature oddity is that it must be turned twice around, through
$720$ degrees rather than $360$, before it returns to itself. A \emph{fermion} is a particle of matter, an
electron or a quark, no two of which can occupy the same state, which is why matter takes up space at all.
Fermions have half-integer spin, so they are described by spinors. The two ideas are linked, and the substrate
delivers that link rather than presupposing it.

\subsection{The spinor coin}

The $24$ local directions of the cell are the $24$-cell, equivalently the $D_4$ root system. Under that symmetry
the $24$ directions decompose as $8v + 8s + 8c$, the $\mathrm{SO}(8)$ vector together with its two spinors,
with triality, the order-three outer symmetry $S_3$, permuting the three eights. The two pieces $8s$ and $8c$ are
genuine half-integer-spin representations. This is the precise sense in which the substrate carries spin at the
base. It is not assumed and not put in by hand. It is a fact about how the local direction set transforms under
its own symmetry group.

\begin{proposition}[The spinor decomposition]
The $24$ directions of the $24$-cell, as the $D_4$ root system, decompose under $\mathrm{SO}(8)$ as
$8v + 8s + 8c$. The summands $8s$ and $8c$ are half-integer-spin representations, so the substrate carries
spinors directly. Triality permutes the three eights and is the seed of three generations.
\end{proposition}

The control case sharpens the claim. The $12$ directions of $\{5,3,4\}$ decompose under the icosahedral
symmetry as $1 + 3 + 3' + 5$, all of integer spin, with no spinor anywhere in the list. A substrate built on
those $12$ directions would carry no fundamental spinor at all. That contrast is one of the reasons the selection
argument lands on $\{3,4,3,4\}$ rather than $\{5,3,4\}$, and it is developed further where the crystallographic
and spinor requirements are set side by side.

\subsection{Matter is fermionic}

A word of care first. A \emph{self}, in the broad sense used later in the paper, is any persistent
self-maintaining pattern, and these range from a single vibe through ever larger patches and networks all the way
up to the entire mesh. Most selves are not particles. What concerns us here is one particular kind of localized
excitation, the topological soliton, which is the matter object the theory identifies with a fundamental fermion.

\begin{definition}[Matter soliton]
A \emph{matter soliton} is a topological knot in the tone field, a \emph{hopfion}, classified by its Hopf charge
in $\pi_3(S^2) = \mathbb{Z}$. The integer charge counts how the tone configuration links with itself, and it
cannot change under any smooth evolution of the field.
\end{definition}

A hopfion of odd Hopf charge picks up a minus sign under a $2\pi$ rotation. This is the Finkelstein-Rubinstein
result, sharpened by the spin-statistics argument of Wilczek and Zee \cite{finkelsteinrubinstein1968,
wilczekzee1983}. A minus sign under a single full turn is the defining mark of a fermion. So the spin-statistics
connection here is geometric. A matter soliton of odd charge is automatically a fermion, with nothing assumed
about it beyond its topology. In the carrying image such a soliton is a braid in the fabric, a knot of fire, and
the topology simply is the weave.

\begin{proposition}[Odd-charge solitons are fermions]
A topological soliton in the tone field with odd Hopf charge transforms with a sign change under a $2\pi$
rotation. By the spin-statistics connection it is a fermion. The matter solitons are therefore fermions as a
consequence of their topology, with no separate statistical postulate.
\end{proposition}

It pays to keep the two ways spin appears in this theory cleanly apart, because they answer different questions.
The first is the \emph{coin spinor}, the $D_4$ pieces $8s$ and $8c$. These are the fundamental fermion fields,
and they genuinely require $\{3,4,3,4\}$, since the $12$-direction alternative carries no spinor at all. The
second is the \emph{topological route}, the identification of a matter soliton with a hopfion and hence with a
fermion. That route is substrate-independent and works even on $\{5,3,4\}$, because it rests only on $\pi_3(S^2)$
and not on the direction set. The coin spinor is the load-bearing one for the fundamental matter fields and for
the $\mathrm{SO}(10)$ gauge story of the next section but one. The topological route is why the matter solitons
come out fermionic. The two agree, but the precise selection claim is that the coin spinor forces $\{3,4,3,4\}$, not that
spin of any kind does.

\subsection{Matter forms and persists}

Why a stabilizer is needed at all is Derrick's theorem. Rescale a soliton by a factor $\lambda$. The exchange
energy scales as $\lambda^{d-2}$ and a Skyrme term scales as $\lambda^{d-4}$. In three dimensions exchange alone
scales as $\lambda$, so the soliton collapses, and a higher-gradient Skyrme term is required to arrest it
\cite{skyrme1962}. In the two-dimensional cusp of $\{5,3,4\}$ exchange is scale-marginal and no stabilizer is
needed, a neat tie to the bulk-cusp dimension split. The live question is therefore whether the rule's collision
supplies a Skyrme term of the stabilizing sign. The Skyrme coefficient appears at the higher-gradient, Burnett
order of the collision's gradient expansion, and the collision is the rule's last free part, fixed only by being
conserving, reversible, local, and lattice-symmetric. So the matter of stability reduces to a definite sign of a
definite quantity.

The sign was measured. The fermion-sea energy of a double-twist periodic texture was evaluated with the kernel
polynomial method, a Chebyshev expansion of the spectral density carried out with a stochastic trace on the GPU.
A single-helix control was subtracted to isolate the Skyrme contribution from the bulk. The isolated Skyrme term
came out positive, that is, stabilizing.

\begin{result}[The Skyrme stabilizer is positive]
The isolated Skyrme coefficient of the collision, measured as the fermion-sea energy of a double-twist periodic
texture with a single-helix control subtracted \cite{pollard2026vibetest}, is positive. Solitons in the tone
field are therefore stable on their own, with no fine-tuning. Matter forms and persists.
\end{result}

The earlier attempts that used localized solitons were confounded by surface-dominated sea energy, and the
double-twist periodic method is the clean measurement that replaces them. Even had the result come out
ambiguous, matter formation would by itself select the collision member with the stabilizing sign. The result
being positive makes that fallback unnecessary.


The matter sector comes from the substrate. The coin carries spinors, the matter solitons are fermions, and the
rule stabilizes matter against collapse. Fermions are not added to the model. They emerge as the stable knots of
the substrate, and the half-integer spin that defines them is already present in the way the cell's $24$ directions
transform.
